ری ایجاد می‌کند؟ (بله — و این نکته‌ای است که حکومت‌های مستبد منکر آن‌اند.)
\end{itemize}
\end{tcolorbox}

\textbf{اخلاق پیروزی:} تاریخ نشان می‌دهد که بسیاری از جنبش‌های آزادی‌خواه 
پس از پیروزی، خود به ستمگران جدید بدل شده‌اند. انقلاب فرانسه به ترور 
ژاکوبنی، انقلاب روسیه به استبداد استالینی، و انقلاب ایران به حکومت 
روحانیون انجامید. \thinker{آلبر کامو} در \textit{انسان طاغی} 
(\lat{L'Homme révolté}، ۱۹۵۱) هشدار داد: «هر انقلابی سرانجام یا سرکوبگر 
می‌شود یا بدعت‌گذار.»\footnote{\lr{Camus, Albert. \textit{The Rebel}. 
1951.}} مسئله این است: چگونه می‌توان از تکرار این چرخهٔ شوم جلوگیری کرد؟

\textbf{اخلاق شکست:} شکست نیز اخلاق خود را دارد. آیا جنبشی که شکست 
می‌خورد، اخلاقاً محکوم است؟ خیر. موفقیت و شکست معیار درستی و نادرستی 
اخلاقی نیست. اما شکست مسئولانه با شکست بی‌مسئولیتانه تفاوت دارد: 
رهبرانی که مردم را بدون آمادگی و برنامه به خیابان می‌کشند و سپس ناپدید 
می‌شوند، اخلاقاً مسئول‌اند — حتی اگر نیتشان خیر بوده باشد.


% ╔══════════════════════════════════════════════════════════════╗
% ║   بخش چهارم: مورد ایران                                   ║
% ╚══════════════════════════════════════════════════════════════╝

\newpage
\section{مورد ایران: تحلیل اخلاقی مبارزه با استبداد دینی}

\subsection{توصیف وضعیت}

جمهوری اسلامی ایران از سال ۱۳۵۷ (۱۹۷۹) تاکنون — یعنی بیش از ۴۷ سال — 
تحت حاکمیت روحانیون شیعی اداره می‌شود. ویژگی‌های ساختاری این حکومت:

\begin{enumerate}[label=\textbf{\arabic*.}]
\item \textbf{تمرکز قدرت در نهاد ولایت فقیه:} رهبر، بالاترین مقام حکومتی 
است که نه انتخابی واقعی است و نه پاسخگو. تمام نهادهای مهم — قوهٔ قضاییه، 
شورای نگهبان، نیروهای مسلح — زیر نظر مستقیم او قرار دارند.
\item \textbf{سرکوب سیستماتیک مخالفان:} از اعدام‌های دههٔ شصت تا سرکوب 
جنبش سبز (۱۳۸۸)، اعتراضات آبان ۱۳۹۸ و جنبش «زن، زندگی، آزادی» (۱۴۰۱)، 
حکومت الگوی ثابتی از سرکوب خشن نشان داده است.
\item \textbf{عدم مشارکت سیاسی واقعی:} انتخابات نمایشی است؛ نامزدها توسط 
شورای نگهبان فیلتر می‌شوند و تنها «خودی‌ها» مجاز به رقابت‌اند.
\item \textbf{کنترل ایدئولوژیک:} حکومت از مذهب به‌عنوان ابزار مشروعیت‌بخشی 
و سرکوب استفاده می‌کند. منتقدان نه‌فقط «ضدانقلاب» بلکه «محارب با خدا» 
(\lat{mohareb}) خوانده می‌شوند.
\item \textbf{فقدان استقلال قوا و عدم حاکمیت قانون:} قوهٔ قضاییه ابزار 
سرکوب است نه مرجع دادخواهی.
\end{enumerate}

\subsection{ارزیابی اخلاقی: آیا شرایط حق براندازی فراهم است؟}

بر اساس معیارهای نظریهٔ جنگ عادلانه و سنت‌های مختلف فلسفی:

\begin{table}[htbp]
\centering
\small
\begin{tabularx}{\linewidth}{|>{\raggedleft\bfseries}p{3.5cm}|X|>{\centering\arraybackslash}p{2cm}|}
\hline
\rowcolor{BgCool}
\textbf{معیار} & \textbf{ارزیابی وضعیت ایران} & \textbf{وضعیت} \\
\hline
علت عادلانه & ستم سیستماتیک مستند و مداوم & \cellcolor{green!15} برقرار \\
\hline
آخرین چاره & اصلاحات، نافرمانی مدنی، اعتراض مسالمت‌آمیز همه آزموده و 
سرکوب شده‌اند & \cellcolor{green!15} برقرار \\
\hline
نیت درست & آزادی‌خواهی و حقوق شهروندی (البته نیاز به تضمین) & 
\cellcolor{yellow!25} نسبی \\
\hline
احتمال موفقیت & بستگی به شرایط، سازمان‌دهی و عوامل بین‌المللی & 
\cellcolor{yellow!25} نامعلوم \\
\hline
تناسب & نیازمند ارزیابی دقیق هزینه-فایده & \cellcolor{yellow!25} نیازمند 
تحلیل \\
\hline
مرجع مشروع & فقدان رهبری واحد و مشروع & \cellcolor{red!15} چالش‌برانگیز \\
\hline
\end{tabularx}
\caption*{ارزیابی شرایط حق براندازی در ایران بر اساس نظریهٔ جنگ عادلانه}
\end{table}

نتیجه آنکه \textit{علت عادلانه} و \textit{آخرین چاره} بودن تقریباً محرز 
است. اما چالش‌های جدی در زمینهٔ \textit{مرجع مشروع}، \textit{احتمال موفقیت} 
و \textit{تناسب} وجود دارد. این بدان معنا نیست که مبارزه ناموجه است — بلکه 
بدان معناست که مبارزه نیازمند سازمان‌دهی، هوشمندی و خویشتنداری بیشتری است.

\subsection{تجربهٔ جنبش‌های اعتراضی ایران: درس‌های اخلاقی}

\subsubsection{جنبش سبز (۱۳۸۸/۲۰۰۹)}
جنبش سبز نمونه‌ای از مبارزهٔ عمدتاً خشونت‌پرهیز بود. نقاط قوت اخلاقی:
\begin{itemize}
\item شعار «رأی من کو؟» مبتنی بر حق مشخص و قابل فهم بود.
\item خشونت‌پرهیزی عمدتاً رعایت شد.
\item مشارکت گستردهٔ اقشار مختلف.
\end{itemize}
نقاط ضعف و درس‌ها:
\begin{itemize}
\item رهبری (موسوی و کروبی) در چارچوب نظام باقی ماند و نتوانست فراتر رود.
\item فقدان برنامهٔ روشن برای مراحل بعد از اعتراض.
\item حکومت با سرکوب شدید جنبش را خاموش کرد و هزینهٔ ادامه‌دادن بسیار بالا 
شد.
\end{itemize}

\subsubsection{اعتراضات آبان ۱۳۹۸ (نوامبر ۲۰۱۹)}
این اعتراضات با قطع اینترنت و کشتار بی‌سابقه (بیش از ۱۵۰۰ نفر به گفتهٔ 
رویترز) سرکوب شد. درس‌های اخلاقی:
\begin{itemize}
\item سرعت و خشونت سرکوب نشان داد که حکومت حاضر به هر هزینه‌ای است.
\item فقدان سازمان‌دهی، معترضان را در برابر سرکوب سازمان‌یافته آسیب‌پذیر کرد.
\item قطع ارتباطات نشان‌دهندهٔ اهمیت زیرساخت‌های ارتباطی مستقل است.
\end{itemize}

\subsubsection{جنبش «زن، زندگی، آزادی» (۱۴۰۱/۲۰۲۲)}
مهم‌ترین و گسترده‌ترین جنبش اعتراضی تاریخ جمهوری اسلامی:
\begin{itemize}
\item شعار «زن، زندگی، آزادی» ماهیتاً \textit{ارزشی و اخلاقی} بود — نه 
صرفاً سیاسی.
\item مشارکت زنان، جوانان و اقلیت‌های قومی بی‌سابقه بود.
\item خشونت‌پرهیزی عمدتاً رعایت شد، هرچند در برخی مناطق (بالاخص 
کردستان و بلوچستان) مقاومت فعال‌تری شکل گرفت.
\item حکومت با ترکیبی از کشتار، بازداشت گسترده، شکنجه و تجاوز سرکوب کرد.
\end{itemize}

\begin{tcolorbox}[quotebox]
\textbf{پرسش اخلاقی تلخ:} وقتی حکومتی حاضر است برای حفظ قدرت، صدها نفر 
را بکشد و هزاران نفر را شکنجه کند — آیا اصرار بر خشونت‌پرهیزی مطلق، 
بار اخلاقی نامتناسبی بر قربانیان نمی‌گذارد؟ آیا «از جان‌گذشتگی» معترضان 
لازمهٔ خشونت‌پرهیزی است یا نتیجهٔ فقدان گزینه‌های دیگر؟
\end{tcolorbox}

\subsection{چالش‌های ویژهٔ مورد ایران}

\subsubsection{تشتت و فقدان رهبری}
یکی از بزرگ‌ترین مشکلات اپوزیسیون ایرانی، تفرق و فقدان انسجام است. 
گروه‌های مختلف — جمهوری‌خواهان، مشروطه‌خواهان، چپ‌ها، اقلیت‌های قومی، 
فعالان مذهبی و سکولار — نه‌تنها بر سر آیندهٔ ایران، بلکه حتی بر سر 
شیوهٔ مبارزه اتفاق نظر ندارند.

از منظر اخلاقی، این تشتت هم مشکل است و هم فرصت:
\begin{itemize}
\item \textbf{مشکل:} بدون حداقلی از اتفاق نظر، هرگونه اقدام جمعی مؤثر 
ناممکن است — و ناتوانی در عمل جمعی خودْ نوعی ناکامی اخلاقی است.
\item \textbf{فرصت:} تکثرگرایی می‌تواند تضمینی باشد علیه تمرکز قدرت پس 
از پیروزی. جنبشی که از ابتدا چندصدایی باشد، کمتر احتمال دارد به 
دیکتاتوری جدید ختم شود.
\end{itemize}

\subsubsection{فقدان حمایت‌های حداقلی}
مبارزه بدون منابع ممکن نیست. اپوزیسیون ایرانی از فقدان جدی در زمینهٔ:
\begin{itemize}
\item حمایت مالی مستقل و پایدار
\item ظرفیت فکری و استراتژیک سازمان‌یافته
\item زمان و فرصت (مردم تحت فشار اقتصادی و امنیتی‌اند)
\item رسانهٔ مستقل و قابل اعتماد
\item ارتباط مؤثر میان داخل و خارج
\end{itemize}
رنج می‌برد. این فقدان‌ها صرفاً «موانع عملی» نیستند — بلکه خود بُعد 
اخلاقی دارند: آیا کسانی که منابع و امکانات دارند (دیاسپورای ایرانی، 
نخبگان خارج از کشور) وظیفهٔ اخلاقی‌ای برای حمایت دارند؟ پاسخ 
اکثر نظریه‌های اخلاقی مثبت است.

\subsubsection{ضعف در اعتراض مدنی: علل ساختاری}
چرا اعتراض مدنی در ایران به نتیجهٔ ماندگار نمی‌رسد؟

\begin{enumerate}
\item \textbf{عدم تقارن قدرت:} حکومت مسلح، سازمان‌یافته و دارای منابع 
مالی بزرگ است. معترضان پراکنده، غیرمسلح و فاقد منابع‌اند.
\item \textbf{اتمیزه‌شدن جامعه:} دهه‌ها سرکوب، بی‌اعتمادی اجتماعی عمیقی 
ایجاد کرده. مردم به یکدیگر اعتماد ندارند و سازمان‌دهی دشوار است.
\item \textbf{خستگی و ناامیدی:} شکست‌های مکرر، «آموخته‌درماندگی» 
(\lat{learned helplessness}) ایجاد کرده.
\item \textbf{هزینهٔ سنگین مشارکت:} زندان، شکنجه، اعدام، از دست دادن 
شغل و تحصیل — هزینهٔ اعتراض در ایران سرسام‌آور است.
\end{enumerate}

\subsection{اخلاق چه حکمی برای مردم ایران دارد؟}

\begin{tcolorbox}[mybox, title={\textbf{حکم اخلاقی ترکیبی}}]
بر اساس تحلیل‌های این مقاله:
\begin{enumerate}[label=\textbf{\arabic*.}]
\item \textbf{مبارزه حق است و حتی وظیفه:} سکوت در برابر ستم سیستماتیک، 
همدستی با آن است. اما این وظیفه مطلق نیست — هر فرد بر اساس توانایی و 
موقعیتش مسئول است (\concept{رابطهٔ توانایی و تعهد}).
\item \textbf{روش‌ها باید اخلاقی باشند:} هدف وسیله را توجیه نمی‌کند — 
اما وسیله هم نباید به بهانهٔ «اخلاق»، هدف (آزادی) را غیرممکن سازد.
\item \textbf{سازمان‌دهی وظیفه‌ای اخلاقی است:} مبارزهٔ بدون سازمان، 
بدون برنامه و بدون ارزیابی ریسک، بی‌مسئولیتی اخلاقی است.
\item \textbf{حمایت از مبارزان وظیفه‌ای مشترک است:} کسانی که خارج از 
دسترس سرکوب هستند — چه ایرانیان خارج از کشور و چه جامعهٔ بین‌المللی — 
مسئولیت اخلاقی حمایت دارند.
\item \textbf{آمادگی برای عدالت انتقالی از همین حالا ضروری است:} 
اندیشیدن به «روز بعد» نه نشانهٔ خوش‌خیالی، بلکه تضمین اخلاقی علیه 
تکرار استبداد است.
\end{enumerate}
\end{tcolorbox}


% ╔══════════════════════════════════════════════════════════════╗
% ║   بخش پنجم: پیچیدگی‌ها، دام‌ها و راه میانه                ║
% ╚══════════════════════════════════════════════════════════════╝

\newpage
\section{پیچیدگی‌ها، دام‌ها و جستجوی راه میانه}

\subsection{دو پرتگاه: ساده‌سازی و پیچیده‌سازی}

یکی از مهم‌ترین بصیرت‌های این مقاله آن است که مسئلهٔ اخلاق براندازی 
همواره میان دو پرتگاه در حرکت است:

\begin{figure}[htbp]
\centering
\begin{adjustbox}{max width=\linewidth}
\begin{tikzpicture}[
  node distance=0.5cm,
  cliff/.style={
    draw, very thick, fill=red!10, rounded corners=3pt,
    minimum width=5cm, minimum height=4cm, align=center,
    font=\small
  },
  bridge/.style={
    draw, very thick, fill=green!10, rounded corners=10pt,
    minimum width=5.5cm, minimum height=3.5cm, align=center,
    font=\small
  },
  arr/.style={-{Stealth[length=6pt]}, thick, color=NeutralDark}
]

% Left cliff
\node[cliff] (left) {%
  \rl{\textbf{\large پرتگاه ساده‌سازی}}\\[6pt]%
  \rl{«هر کاری برای آزادی مجاز است»}\\[3pt]%
  \rl{«دشمن = شیطان مطلق»}\\[3pt]%
  \rl{«اخلاق = ضعف»}\\[3pt]%
  \rl{«فقط نتیجه مهم است»}\\[6pt]%
  \rl{\color{AccentRed}\textbf{نتیجه: بازتولید استبداد}}%
};

% Right cliff
\node[cliff, right=5.5cm of left] (right) {%
  \rl{\textbf{\large پرتگاه پیچیده‌سازی}}\\[6pt]%
  \rl{«شرایط هنوز فراهم نیست»}\\[3pt]%
  \rl{«هیچ عملی کامل نیست»}\\[3pt]%
  \rl{«اخلاق = بی‌عملی»}\\[3pt]%
  \rl{«همه‌چیز خاکستری است»}\\[6pt]%
  \rl{\color{AccentRed}\textbf{نتیجه: تداوم استبداد}}%
};

% Bridge in middle
\node[bridge] at ($(left.east)!0.5!(right.west)+(0,4.5)$) (bridge) {%
  \rl{\textbf{\large\color{AccentGreen} راه میانه}}\\[6pt]%
  \rl{عمل سنجیده + نقد مداوم}\\[3pt]%
  \rl{اصول روشن + انعطاف عملی}\\[3pt]%
  \rl{اخلاق = قدرت و هوشمندی}\\[3pt]%
  \rl{هم فرایند مهم و هم نتیجه}%
};

\draw[arr] (bridge.south west) -- (left.north east);
\draw[arr] (bridge.south east) -- (right.north west);
\draw[arr, dashed, color=AccentRed] (left.north) -- ++(0,1) 
  node[above, font=\tiny, text=AccentRed] {\rl{سقوط}};
\draw[arr, dashed, color=AccentRed] (right.north) -- ++(0,1) 
  node[above, font=\tiny, text=AccentRed] {\rl{فلج}};

\end{tikzpicture}
\end{adjustbox}
\caption*{دو پرتگاه و راه میانه}
\end{figure}

\subsubsection{خطرات ساده‌سازی افراطی}
\begin{enumerate}
\item \textbf{شیطانی‌سازی مطلق دشمن:} وقتی حکومت به‌کلی غیرانسانی تصویر 
شود، هر خشونتی علیه هر کسی که با حکومت مرتبط باشد موجه جلوه می‌کند. 
این راه به تروریسم، قتل‌عام و بی‌عدالتی جدید ختم می‌شود.
\item \textbf{نفی هرگونه محدودیت اخلاقی:} «در جنگ با شیطان، هر وسیله‌ای 
مقدس است» — این منطق دقیقاً همان منطقی است که ستمگران هم به کار می‌برند.
\item \textbf{شخصی‌شدن قدرت:} وقتی مبارزه «ساده» تعریف شود، رهبران 
کاریزماتیک جای اندیشهٔ جمعی را می‌گیرند و خطر دیکتاتوری جدید بالا می‌رود.
\item \textbf{از دست دادن متحدان میانه‌رو:} ساده‌سازی، طیف خاکستری جامعه 
— شامل افرادی که نه با حکومت‌اند و نه با اپوزیسیون — را از جنبش دور می‌کند.
\end{enumerate}

\subsubsection{خطرات پیچیده‌سازی افراطی}
\begin{enumerate}
\item \textbf{فلج تحلیلی (\lat{analysis paralysis}):} «هنوز شرایط فراهم 
نیست»، «باید بیشتر بررسی کنیم»، «بیایید اول اختلافاتمان را حل کنیم» — 
این عبارات هرچند معقول به نظر می‌رسند، می‌توانند بهانه‌ای برای بی‌عملی 
ابدی شوند.
\item \textbf{نسبی‌گرایی اخلاقی:} «همه‌چیز پیچیده است» می‌تواند به 
«هیچ‌چیز واقعاً درست یا غلط نیست» تبدیل شود — و این مستقیماً به نفع 
وضع موجود تمام می‌شود.
\item \textbf{انتقال بار اخلاقی به قربانی:} وقتی از معترض خواسته می‌شود 
«اخلاقی» باشد بدون آنکه هیچ انتظار مشابهی از حکومت وجود داشته باشد، 
اخلاق به ابزار سرکوب تبدیل شده.
\item \textbf{بهانه‌تراشی برای بی‌عملی نخبگان:} روشنفکرانی که با تحلیل‌های 
بی‌پایان از عمل شانه خالی می‌کنند، از پیچیدگی به‌عنوان سپر استفاده می‌کنند.
\end{enumerate}

\subsection{اخلاقِ مانع و اخلاقِ یاری‌رسان: نمونه‌های تاریخی}

\begin{tcolorbox}[successbox, title={\textbf{نمونه‌های اخلاق یاری‌رسان}}]
\begin{enumerate}
\item \textbf{جنبش ضدآپارتاید آفریقای جنوبی:} \thinker{نلسون ماندلا} 
ابتدا مسیر خشونت‌پرهیز را طی کرد، سپس به مبارزهٔ مسلحانهٔ محدود 
(خرابکاری در زیرساخت‌ها، نه حمله به غیرنظامیان) روی آورد، و پس از ۲۷ 
سال زندان، عدالت انتقالی مبتنی بر آشتی (\lat{Truth and Reconciliation 
Commission}) را پیشه کرد. او توانست بدون ساده‌سازی (نه همهٔ سفیدپوستان 
دشمن بودند) و بدون پیچیده‌سازی فلج‌کننده (نباید منتظر شرایط ایدئال ماند) 
راه میانه‌ای پیدا کند.

\item \textbf{جنبش حقوق مدنی آمریکا:} \thinker{مارتین لوتر کینگ} با ترکیب 
اصول اخلاقی روشن (کرامت انسانی، خشونت‌پرهیزی) و عمل‌گرایی استراتژیک 
(انتخاب دقیق زمان و مکان اعتراض، جلب توجه رسانه‌ها) توانست جنبشی بسازد 
که هم اخلاقی بود و هم مؤثر.

\item \textbf{انقلاب مخملی چکسلواکی (۱۹۸۹):} \thinker{واتسلاو هاول} 
نویسنده‌ای بود که با مفهوم «زیستن در حقیقت» (\lat{Living in Truth}) 
نشان داد که حتی در دل توتالیتاریسم، پایبندی به اخلاق و حقیقت می‌تواند 
قدرت سیاسی داشته باشد. موفقیت انقلاب مخملی مدیون دهه‌ها مقاومت فرهنگی 
و اخلاقی صبورانه بود.

\item \textbf{جنبش همبستگی لهستان (\lat{Solidarność}):} 
\thinker{لخ واونسا} و جنبش همبستگی نشان دادند که ترکیب مبارزهٔ کارگری 
سازمان‌یافته، حمایت کلیسا و خشونت‌پرهیزی می‌تواند حتی یک حکومت 
کمونیستی تحت حمایت شوروی را سرنگون کند.
\end{enumerate}
\end{tcolorbox}

\begin{tcolorbox}[enemybox, title={\textbf{نمونه‌های اخلاقِ مانع آزادی}}]
\begin{enumerate}
\item \textbf{سنت «اطاعت از حاکم» در تفکر اسلامی محافظه‌کار:} حدیث 
«اُصبِروا» (صبر کنید حتی اگر حاکم شما را بزند و مالتان را بگیرد) 
قرن‌ها برای توجیه اطاعت از حاکمان ستمگر به کار رفته. این «اخلاق» 
در خدمت قدرت بود، نه عدالت.

\item \textbf{«نظم‌گرایی» در برابر جنبش حقوق مدنی:} کینگ در \textit{نامه 
از زندان بیرمنگام} از «سفیدپوست میانه‌رو» (\lat{white moderate}) نوشت 
که برایش «نظم» مهم‌تر از «عدالت» بود: «متأسفانه بزرگ‌ترین سنگ راهمان نه 
عضو کلان‌ها، بلکه سفیدپوست میانه‌روی است که نظم را بر عدالت ترجیح 
می‌دهد.»\footnote{\lr{King Jr., Martin Luther. ``Letter from Birmingham 
Jail.'' 1963.}}

\item \textbf{«واقع‌گرایی» استعماری:} استعمارگران همواره از «عدم آمادگی 
بومیان برای خودحکومتی» سخن می‌گفتند — نوعی 
پیچیده‌سازی اخلاقی‌نما برای تداوم سلطه.

\item \textbf{«آرامش‌طلبی» در ایران:} پس از هر سرکوبی، صداهایی بلند 
می‌شوند که می‌گویند «حالا وقتش نیست»، «تنش نکنید»، «به فکر زندگیتان 
باشید». این صداها گاهی از دلسوزی واقعی برمی‌آیند — اما عملاً در خدمت 
تداوم وضع موجود‌اند.

\item \textbf{اخلاق «کویتیسم» (\lat{Quietism}) در تشیع سنتی:} 
بسیاری از مراجع تقلید شیعه با استناد به غیبت امام زمان، هرگونه قیام 
سیاسی پیش از ظهور را ناروا دانسته‌اند. این دیدگاه عملاً به پذیرش هر 
حکومتی — هرقدر ستمگر — منجر می‌شود.
\end{enumerate}
\end{tcolorbox}

\subsection{رابطهٔ توانایی، تعهد و عمل}

\begin{tcolorbox}[defbox, title={\textbf{اصل تناسب مسئولیت با توانایی}}]
یکی از اصول مهم اخلاقی که در بحث براندازی اغلب نادیده گرفته می‌شود:
\textbf{مسئولیت اخلاقی هر فرد متناسب با توانایی اوست.}

\begin{itemize}
\item از مادری که تنها نان‌آور خانواده‌اش است، نمی‌توان انتظار حضور در 
خط مقدم داشت.
\item از روشنفکری که در امنیت خارج از کشور زندگی می‌کند، می‌توان و باید 
انتظار بیشتری داشت.
\item از سازمان‌های بین‌المللی حقوق بشر، انتظار حمایت مؤثر — نه صرفاً 
بیانیه — می‌رود.
\item از نظامیان و نیروهای امنیتی حکومت، انتظار سرپیچی از دستور 
سرکوب — دست‌کم خودداری از مشارکت فعال — اخلاقاً موجه است.
\end{itemize}
\end{tcolorbox}

\subsection{مسئلهٔ حقوق مشاع و زیاده‌خواهی}

\begin{tcolorbox}[warnbox, title={\textbf{چالش در حقوق مشاع}}]
وقتی گروهی در امور مشاع (سیاست، منابع ملی، فضای عمومی) زیاده‌خواهی و 
زورگویی می‌کند — با فرض اینکه قدرت بیشتری دارد — چه حقوقی برای واکنش 
وجود دارد؟

نکتهٔ مهم: \textbf{وقتی حقوقی زایل شده، برای از دست دهندگانش اهمیت 
روزافزون می‌یابد.} این یک واقعیت روان‌شناختی است: آزادی‌ای که نداریم، 
بیش از آزادی‌ای که داریم ارزش پیدا می‌کند. حکومت مستبد با سلب حقوق مردم، 
ناخواسته اهمیت آن حقوق را در ذهن مردم صدچندان می‌کند — و این خود سوختِ 
مبارزه می‌شود.
\end{tcolorbox}

\subsection{الاهیات مبارزه و آزادی‌خواهی}

نقش دین و الاهیات در مبارزات آزادی‌خواهانه دوگانه بوده است: هم ابزار 
سرکوب و هم منبع الهام مقاومت. 

\begin{table}[htbp]
\centering
\small
\begin{tabularx}{\linewidth}{|>{\raggedleft\bfseries}p{3cm}|X|X|}
\hline
\rowcolor{BgCool}
\textbf{سطح} & \textbf{کاربرد سرکوبگرانه} & \textbf{کاربرد رهایی‌بخش} \\
\hline
الاهیات سیاسی & حق الاهی حاکم، ولایت فقیه & الاهیات رهایی‌بخش، عدالت 
الاهی \\
\hline
ایده‌ها و آرمان‌ها & تقدیرگرایی، صبر در برابر ظلم & عدالت‌خواهی، کرامت 
انسانی \\
\hline
کنش سیاسی & فتوای سرکوب، تکفیر مخالفان & بسیج مردم، شبکه‌های اجتماعی 
مسجد \\
\hline
اصلاح و نقد & سانسور، ممنوعیت اجتهاد & نواندیشی دینی، اصلاح از درون \\
\hline
\end{tabularx}
\caption*{دوگانگی نقش الاهیات در مبارزه و سرکوب}
\end{table}

\thinker{الاهیات رهایی‌بخش} (\lat{Liberation Theology}) — که در آمریکای 
لاتین توسط \thinker{گوستاوو گوتیرز} (\lat{Gustavo Gutiérrez}) و دیگران 
توسعه یافت — تلاش کرد مسیحیت را از ابزار قدرت به منبع مقاومت تبدیل کند. 
در ایران، \thinker{علی شریعتی} و در دورهٔ اخیر \thinker{محمد مجتهد شبستری} 
و \thinker{عبدالکریم سروش} تلاش‌هایی مشابه در بستر اسلام صورت داده‌اند — 
هرچند این تلاش‌ها با محدودیت‌ها و انتقادات جدی نیز مواجه بوده‌اند.

\subsection{دیدن ضعف‌ها و مشکلات خودی}

\begin{tcolorbox}[warnbox, title={\textbf{آینه‌ای ناخوشایند اما ضروری}}]
اخلاق مبارزه ایجاب می‌کند که مبارز نه‌فقط ستم حکومت، بلکه ضعف‌ها و 
مشکلات خودش و جنبشش را نیز ببیند:
\begin{itemize}
\item \textbf{ترس‌ها:} پذیرفتن ترس، نه سرکوب آن. شجاعت غیاب ترس نیست، 
بلکه عمل‌کردن با وجود ترس است.
\item \textbf{سستی‌ها:} بی‌عملی، حتی وقتی توجیه اخلاقی‌نما دارد، باید 
صادقانه بازشناسی شود.
\item \textbf{تمایلات سلطه‌جویانهٔ درونی:} بسیاری از مبارزان در درون خود 
الگوهای ستمگرانه‌ای حمل می‌کنند — سلطه‌جویی مردانه، قومیت‌گرایی، 
بزرگ‌منشی طبقاتی. مبارزه‌ای که این‌ها را بازشناسی نکند، محکوم به 
تکرارشان است.
\item \textbf{اختلافات درون‌گروهی:} نه ماست‌مالی کردن و نه به بهانهٔ 
«وحدت» سرکوب کردن صداهای مخالف — بلکه مدیریت دموکراتیک اختلاف.
\end{itemize}
\end{tcolorbox}


% ╔══════════════════════════════════════════════════════════════╗
% ║         بخش ششم: نتیجه‌گیری                                ║
% ╚══════════════════════════════════════════════════════════════╝

\newpage
\section{نتیجه‌گیری: به سوی اخلاقی توانمندساز}

\subsection{خلاصهٔ یافته‌ها}

این مقاله نشان داد که:

\begin{enumerate}[label=\textbf{\arabic*.}]
\item \textbf{حق مقاومت و براندازی} در شرایط ستم سیستماتیک، در تمامی 
سنت‌های فکری بزرگ به رسمیت شناخته شده — اما همواره مشروط و محدود بوده 
است.

\item \textbf{محدودیت‌های اخلاقی} نه ضعف، بلکه منبع قدرت‌اند: جنبشی که 
اصول اخلاقی‌اش را حفظ کند، مشروعیت بیشتری دارد، حمایت گسترده‌تری جلب 
می‌کند و احتمال بیشتری دارد که به نظم عادلانه‌تری ختم شود.

\item \textbf{خطر اصلی} نه در داشتن اصول اخلاقی، بلکه در سوءتعبیر آنها 
به‌عنوان حکم بی‌عملی است. اخلاقی که مبارز را فلج کند، اخلاق نیست — 
ابزار سرکوب است.

\item \textbf{راه میانه} عبارت است از: عملِ سنجیده با نقد مداوم، اصولِ 
روشن با انعطاف عملی، و پایبندی به کرامت انسانی حتی در قلب نبرد.

\item \textbf{مورد ایران} نشان‌دهندهٔ همهٔ پیچیدگی‌های مسئله است: 
مردمی با علت عادلانه اما بدون ابزار کافی، با دشمنی بی‌رحم اما نه 
شکست‌ناپذیر، و با نیاز مبرم به سازمان‌دهی، حمایت و اندیشهٔ اخلاقی 
تواناساز.
\end{enumerate}

\subsection{فراخوان: اخلاقی که فلج نمی‌کند}

\begin{tcolorbox}[mybox, title={\textbf{بیانیهٔ پایانی: اخلاق به‌مثابه 
قدرت}}]
اخلاق براندازی نه مجوز بی‌قیدوشرط خشونت است و نه فرمان بی‌عملی. 
اخلاق، وقتی درست فهمیده شود، \textit{قدرت} است:

\begin{itemize}
\item قدرتِ تمییز دادن میان عدالت و انتقام.
\item قدرتِ حفظ انسانیت در دل نبرد.
\item قدرتِ سازمان‌دهی بر اساس اعتماد و نه ترس.
\item قدرتِ جلب حمایت جهانی از طریق مشروعیت.
\item قدرتِ ساختن آینده‌ای متفاوت — نه تکرار گذشته با بازیگران جدید.
\end{itemize}

مردمانی که تحت ستم می‌زیند، حق دارند مبارزه کنند. و مبارزه‌ای که 
اخلاقی باشد، نه‌تنها ممکن و مطلوب، بلکه \textbf{مؤثرتر} نیز هست. 
تاریخ شاهد است: جنبش‌هایی که اصولشان را حفظ کردند، نه‌فقط پیروز شدند، 
بلکه پیروزی‌هایشان پایدارتر بود.

\vspace{8pt}
\begin{center}
\textit{«در دشمنی با شیطان، به شیطان بدل نشدن — این بزرگ‌ترین پیروزی 
اخلاقی مبارز است.»}
\end{center}
\end{tcolorbox}


% ╔══════════════════════════════════════════════════════════════╗
% ║                    پیوست‌ها                                ║
% ╚══════════════════════════════════════════════════════════════╝

\newpage
\appendix

\section*{پیوست‌ها}
\addcontentsline{toc}{section}{پیوست‌ها}

% ============ پیوست الف: مرور تاریخی و موج‌شناسی ============

\subsection*{پیوست الف: مرور تاریخی و موج‌شناسی جنبش‌های مقاومت}
\addcontentsline{toc}{subsection}{پیوست الف: مرور تاریخی و موج‌شناسی}

\subsubsection*{موج اول: عصیان‌های پیشامدرن (تا قرن هجدهم)}

عصیان علیه حاکمان ستمگر پیشینه‌ای به قدمت خودِ تمدن دارد. از شورش 
بردگان اسپارتاکوس در رم (۷۳-۷۱ ق.م.) تا جنبش‌های هزاره‌گرایانهٔ 
(\lat{millenarian}) قرون وسطی، مردمان تحت ستم همواره راه‌هایی برای 
مقاومت یافته‌اند. ویژگی‌های مشترک:

\begin{itemize}
\item فقدان ایدئولوژی سیاسی مدرن — عصیان بیشتر از روی ناچاری و 
ستم غیرقابل‌تحمل بود تا برنامهٔ سیاسی.
\item رهبری کاریزماتیک — شخصیت‌محور نه نهادمحور.
\item خشونت گسترده و بی‌قاعده — مفهوم «قوانین جنگ» هنوز شکل نگرفته بود.
\item سرکوب بی‌رحمانه و معمولاً شکست.
\end{itemize}

\textbf{نمونه‌ها:}
\begin{itemize}
\item شورش اسپارتاکوس (۷۳ ق.م.)
\item جنبش‌های خوارج در اسلام اولیه (قرن اول هجری)
\item شورش دهقانان آلمان (۱۵۲۴-۱۵۲۵) — با رهبری فکری \thinker{توماس 
مونتسر}
\item قیام پوگاچف در روسیه (۱۷۷۳-۱۷۷۵)
\end{itemize}

\subsubsection*{موج دوم: انقلاب‌های مدرن (قرن هجدهم تا نوزدهم)}

انقلاب‌های آمریکا (۱۷۷۶)، فرانسه (۱۷۸۹) و جنبش‌های ملی‌گرایانهٔ قرن 
نوزدهم. ویژگی‌ها:

\begin{itemize}
\item ایدئولوژی سیاسی صریح — حقوق طبیعی، حاکمیت مردم.
\item اسناد بنیادین — اعلامیهٔ استقلال، اعلامیهٔ حقوق بشر و شهروند.
\item تلاش برای نهادسازی پس از انقلاب.
\item خشونت عمدتاً در چارچوب جنگ نظامی.
\end{itemize}

\textbf{نمونه‌ها:}
\begin{itemize}
\item انقلاب آمریکا (۱۷۷۶)
\item انقلاب فرانسه (۱۷۸۹) — با فازهای مختلف از آزادی‌خواهی تا ترور
\item انقلاب‌های ۱۸۴۸ اروپا (بهار ملت‌ها)
\item جنبش مشروطه ایران (۱۲۸۵-۱۲۹۰ شمسی / ۱۹۰۵-۱۹۱۱)
\end{itemize}

\subsubsection*{موج سوم: انقلاب‌های ضداستعماری و رهایی‌بخش (قرن بیستم)}

ویژگی‌ها:
\begin{itemize}
\item ترکیب ناسیونالیسم با ایدئولوژی‌های مدرن (مارکسیسم، اسلام‌گرایی).
\item جنبش‌های چریکی و پارتیزانی.
\item نقش بین‌المللی برجسته — جنگ سرد، سازمان ملل.
\item بحث جدی دربارهٔ اخلاق خشونت (فانون، سارتر).
\end{itemize}

\textbf{نمونه‌ها:}
\begin{itemize}
\item استقلال هند (۱۹۴۷) — الگوی خشونت‌پرهیزی گاندی
\item انقلاب الجزایر (۱۹۵۴-۱۹۶۲) — الگوی مبارزهٔ مسلحانهٔ فانون
\item انقلاب کوبا (۱۹۵۹) — چه گوارا و مدل چریکی
\item جنبش‌های آزادی‌بخش آفریقا (۱۹۶۰-۱۹۹۰)
\item انقلاب ایران (۱۹۷۹/۱۳۵۷) — ترکیب منحصربه‌فرد اسلام‌گرایی و بسیج 
تودهٔ مردم
\end{itemize}

\subsubsection*{موج چهارم: انقلاب‌های دموکراتیک و مدنی (۱۹۸۹ به بعد)}

ویژگی‌ها:
\begin{itemize}
\item خشونت‌پرهیزی به‌عنوان استراتژی اصلی.
\item نقش رسانه‌ها و افکار عمومی.
\item مفهوم «جامعهٔ مدنی» به‌عنوان بازیگر اصلی.
\item الگوی «انقلاب مخملی».
\end{itemize}

\textbf{نمونه‌ها:}
\begin{itemize}
\item سقوط دیوار برلین و انقلاب‌های ۱۹۸۹ در اروپای شرقی
\item انقلاب مخملی چکسلواکی (۱۹۸۹)
\item فروپاشی آپارتاید (۱۹۹۴)
\item انقلاب‌های رنگی: گرجستان (۲۰۰۳)، اوکراین (۲۰۰۴)، قرقیزستان (۲۰۰۵)
\end{itemize}

\subsubsection*{موج پنجم: جنبش‌های شبکه‌ای عصر دیجیتال (۲۰۱۰ به بعد)}

ویژگی‌ها:
\begin{itemize}
\item نقش محوری شبکه‌های اجتماعی و اینترنت.
\item ساختار افقی و بدون رهبری متمرکز.
\item سرعت بالای بسیج اما شکنندگی سازمانی.
\item پاسخ حکومت‌ها: قطع اینترنت، نظارت دیجیتال.
\end{itemize}

\textbf{نمونه‌ها:}
\begin{itemize}
\item بهار عربی (۲۰۱۰-۲۰۱۲) — موفقیت در تونس، شکست در مصر، فاجعه در سوریه
\item جنبش اشغال وال‌استریت (۲۰۱۱)
\item جنبش چتر هنگ‌کنگ (۲۰۱۴ و ۲۰۱۹)
\item جنبش «زن، زندگی، آزادی» ایران (۲۰۲۲)
\item اعتراضات سراسری سودان (۲۰۱۹ و ۲۰۲۱)
\end{itemize}

\subsubsection*{جدول تطبیقی موج‌ها}

\begin{table}[htbp]
\centering
\small
\begin{adjustbox}{max width=\linewidth}
\begin{tabularx}{1.1\linewidth}{|>{\raggedleft\bfseries\small}p{1.8cm}|
  >{\small}p{2cm}|>{\small}p{2.2cm}|>{\small}p{2.2cm}|
  >{\small}p{2cm}|>{\small}p{2.5cm}|}
\hline
\rowcolor{BgCool}
\textbf{ویژگی} & \textbf{موج ۱} & \textbf{موج ۲} & \textbf{موج ۳} & 
\textbf{موج ۴} & \textbf{موج ۵} \\
\hline
دوره & پیشامدرن & ۱۷۷۰-۱۸۵۰ & ۱۹۴۰-۱۹۸۰ & ۱۹۸۹-۲۰۰۵ & ۲۰۱۰ تاکنون \\
\hline
ایدئولوژی & مذهبی/عرفی & حقوق طبیعی & مارکسیسم/ملی‌گرایی & دموکراسی/حقوق 
بشر & شبکه‌ای/پساایدئولوژیک \\
\hline
روش غالب & خشونت بی‌قاعده & جنگ انقلابی & چریکی/توده‌ای & خشونت‌پرهیز & 
ترکیبی/دیجیتال \\
\hline
رهبری & کاریزماتیک & نخبه‌ای & حزبی & مدنی & افقی/بدون رهبر \\
\hline
نتیجهٔ غالب & شکست & متنوع & استقلال اما استبداد جدید & دموکراسی شکننده & 
نامعلوم \\
\hline
\end{tabularx}
\end{adjustbox}
\caption*{مقایسهٔ تطبیقی پ

ادامهٔ مقاله: